\section{Discussion}

Across four adapted vignette experiments, LLM treatment sensitivity varies by domain, model, and surrounding system instructions. Under the historical no-stimulus condition, Hassner treatment contrasts have the same direction as in the human study, whereas the Savary public-recognition contrast has the opposite sign. Bolsen shows partial directional correspondence but often different magnitudes. Corbacho produces uniformly negative answers under the neutral-realism reference N0; the corresponding human direct-response contrast is small and imprecisely estimated, limiting claims about replication failure.

The original A/C stimuli move Hassner and Savary answers away from the normatively approved response direction. This pattern is not general across the additional cases: relative to N0, A raises Bolsen energy-conservation willingness in GLM but lowers it in Kimi. In Corbacho, A/C change Kimi's affirmative and refusal shares substantially, and the institutional-inefficiency placebo elicits more affirmative answers than the corruption-information treatment. The new N$-$/N$+$ shifts are generally much smaller. These patterns establish prompt sensitivity, not greater fidelity to humans or a uniform reduction in social desirability bias. N0 includes a realism instruction and is not equivalent to historical no stimulus.

The historical H--M comparisons suggest an association between framing valence and responses, but the incomplete prompt factorial does not isolate valence causally. Reasoning annotations also vary in content and local argumentative direction across prompts, yet jointly generated traces and answers cannot establish mediation. Exploratory hidden-state differences in a separate local Qwen3-8B likewise cannot explain the API models' responses. LLM simulations may help stress-test experimental wording, but comparisons against human treatment benchmarks remain necessary.

\section{Scope and Limitations}

The four cases cover selected socially sensitive vignettes, not representative human survey populations. Model coverage is unequal: the historical H--M family is limited to GLM and Kimi, and the principal new-experiment comparisons likewise use those two models. Repeated low-temperature generations are not independent human respondents. Corbacho's three nominal answers require separate treatment of refusal; its human privacy-adjusted estimate is not directly equivalent to our direct-response outcome. The N0 reference and historical no-stimulus baseline differ, as do some original vignette procedures.

The historical reasoning codebook was checked against only 27 expert-coded traces, with highly uneven agreement across categories. Newly planned Bolsen/Corbacho reasoning annotation is not yet part of these results. Reasoning-answer associations are predictive, not causal. Hidden-state measurements concern 98 runs of a different, local Qwen3-8B checkpoint; layer-wise geometry and probing are diagnostic rather than mechanistic proof. Historical H--M prompt wording has not been independently authenticated verbatim. Finally, the historical numerical reconciliation uses an available master-table copy and does not establish identity with the original notebook input.

\section{Conclusion}

System prompts materially affect synthetic experimental answers, but the direction and implications of those effects depend on task and model. Agreement with some human treatment contrasts coexists with reversals, attenuated effects, response saturation, and placebo sensitivity. Historical traces and separate-model representations provide complementary correlational evidence, not a validated causal pathway. Consequently, moving answers away from socially approved positions must not be equated with improved human realism. LLMs are potentially useful for experimental stress-testing, but not as unvalidated substitutes for human respondents.
